\section{Title, outline, and background}
\label{sec:front}

\subsection{Title page and outline (slides 1--2)}

\begin{whybox}
The title page names the project, authors, supervisor, and department.
The outline shows the five big blocks of the talk: Background, Problem/Objectives,
Methodology (including mathematics), Results, Impact.
\end{whybox}

\begin{saybox}
``Today we present a hospital chatbot that reads an English chief complaint and returns
a SATS colour plus a pathway card for \examplehospital{}. The talk moves from the hospital
problem, through the mathematics of BioBERT triage, to holdout results and anticipated impact.''
\end{saybox}

\paragraph{Core idea of the project.}
Patients often arrive with \emph{words only}. Nurses later measure vitals (TEWS).
Our system acts at the text-first moment: free-text complaint \(X\) \(\to\) acuity probabilities
\(\hat{\mathbf{y}}\) \(\to\) colour \(C_{\mathrm{NLP}}\) \(\to\) pathway card \(P(C)\).
Nurses keep TEWS and the right to override.

\subsection{Hospital overcrowding photo (slide 3)}

\begin{whybox}
A visual of overcrowded wards grounds the problem in Ghanaian hospital reality before any equations.
Source attribution: GhanaFact fact-check article on Korle-Bu overcrowding imagery.
\end{whybox}

\begin{saybox}
``This is the world our system sits in: too many patients, limited beds and staff.
Triage has to decide who is seen first and where they go next.''
\end{saybox}

\subsection{Hospital intake under pressure (slide 4)}

\begin{whybox}
Defines the operational setting: crowded queues, SATS at teaching hospitals including
\examplehospital{}, and \emph{acuity discontinuity}---colour assigned at the first desk may not travel with the patient.
\end{whybox}

\textbf{Acuity discontinuity} means: urgency is decided once at intake, then often forgotten when the patient moves to imaging, labs, or wards. Colour and countdown do not reliably steer later decisions \citep{Rominski2014,kathsats}.

\begin{saybox}
``Even when SATS is used, urgency can be lost after the first desk. That is the discontinuity we want a digital colour-plus-pathway record to reduce.''
\end{saybox}

\subsection{SATS colour bands (slide 5)}

\begin{whybox}
Everyone in the room must share the same colour language before you talk about NLP.
\end{whybox}

\begin{definition}[SATS colours]
The South African Triage Scale uses five colours \citep{satsmanual}:
\begin{itemize}[nosep]
  \item \textbf{Red:} immediate / resuscitation stream.
  \item \textbf{Orange:} very urgent (about \(<10\) minutes).
  \item \textbf{Yellow:} urgent (about \(<60\) minutes).
  \item \textbf{Green:} routine / outpatient stream.
  \item \textbf{Blue:} deceased-on-arrival dignity protocol (not ``low urgency'').
\end{itemize}
\end{definition}

Our chatbot maps predicted acuity levels to Red/Orange/Yellow/Green for living patients.
Blue is a dignity protocol, not an NLP output class in this prototype.

\subsection{Text first, then vitals (slide 6)}

\begin{whybox}
Separates the two clocks of triage: words arrive before thermometers.
\end{whybox}

\begin{itemize}
  \item \textbf{Chief complaint:} primary reason for visit, in the patient's (or attendant's) own words.
  \item \textbf{TEWS:} nurse score from HR, RR, mobility, temperature, AVPU, trauma \(\to\) colour from vitals.
  \item \textbf{This project:} English text \(\to\) colour + pathway \emph{before} vitals; TEWS stays in parallel.
\end{itemize}

\begin{saybox}
``We do not replace the nurse. We give a first acuity signal from English text while TEWS remains the vital-sign standard.''
\end{saybox}

\subsection{Literature review (slide 7)}

\begin{whybox}
Shows you know the field: overcrowding/SATS, TEWS limits, BioBERT lineage, NLP triage gaps.
\end{whybox}

\paragraph{Four points to master.}
\begin{enumerate}
  \item LMIC overcrowding and SATS with waiting times \(T_{\max}\) \citep{kathsats,satsmanual}; text-first arrival still under-formalised.
  \item TEWS is objective and does \textbf{not} read language \citep{satsmanual}.
  \item Self-attention \(\to\) BERT \(\to\) BioBERT on PubMed/PMC \citep{vaswani2017attention,lee2020biobert}.
  \item NLP at ED triage is active but mostly high-income; local SATS+pathway math is thin \citep{stewart2023nlp,gligorijevic2018deep}.
\end{enumerate}

\begin{saybox}
``The literature gives us BioBERT and triage NLP abroad. What is missing is an examinable map from Ghanaian intake text to SATS colour and a local pathway card---that is our gap.''
\end{saybox}
